\chapter*{ЗАКЛЮЧЕНИЕ}
\addcontentsline{toc}{chapter}{ЗАКЛЮЧЕНИЕ}

В работе предложен и реализован подход к генерации системно-динамических моделей, в котором большая языковая модель используется для интерпретации описания, а критические свойства результата обеспечиваются типизированным представлением и детерминированными проверками. Это позволяет отделить вероятностную часть генерации от вычисления производных, проверки units, топологии и исполняемости.

Разработаны канонический ModelDocument, ограниченный AST, типизированные компоненты, построитель явных ОДУ и ограниченный repair. Все представления --- CLD, SFD, Python, Vensim, XMILE, симуляция и анализ --- порождаются из единого документа. Пользовательский произвольный код не выполняется.

Созданы generation core, backend и frontend, поддерживающие CLI, REST, email/OTP, контекстные беседы, версии моделей, сценарии, анализ чувствительности и калибровку. Реализация проверена 319 core-тестами, 20 backend и 12 frontend-тестами, а также двумя браузерными E2E-сценариями. Дополнительный изолированный REST E2E охватил 46 HTTP-шагов и 28 проверок: OTP/refresh, порядок сообщений, идемпотентность, версии, симуляцию, exports и ownership.

На SDXorg импортированы и симулированы 63 модели, из которых 62 совпали с эталоном при неизменных допусках. Это составляет 98,41\% внутри поддерживаемого подмножества и 41,89\% всех 148 reference-моделей. Реализованы массивы, поддиапазоны, `:EXCEPT:` и bounded macros без выполнения произвольного кода. Rössler продемонстрировал важность различия reference precision.

Более ранняя acceptance-серия локальной Qwen3-8B-Q4\_K\_M показала 72,86\% executable и semantic и 71,43\% behavior на 70 генерациях. В новом controlled blind ablation на 216 случаях полный pipeline этой модели через llama.cpp достиг 94,44\% executable, 87,50\% semantic и 84,72\% behavior. Direct дал 2,78\% по всем метрикам, typed pipeline без repair --- 73,61\% executable, 69,44\% semantic и 68,06\% behavior. Таким образом, вклад типизированного представления и repair подтверждён парными измерениями, но behavior остаётся на 0,28 п.п. ниже инженерного порога 85\%.

Для локальной 8B repair увеличил behavior на 16,67 п.п. (95\% cluster CI [8,33; 26,39] п.п., Holm $p=0{,}000488$). На Ministral-8B-2512 полный pipeline ранее достиг 97,22\% executable и 90,28\% semantic/behavior, однако эффект repair на semantic/behavior был нулевым ($p=1,0$). Следовательно, вклад repair зависит от модели и конфигурации. Результаты не переносятся на другие checkpoints, providers или произвольные предметные области.

Основные ограничения связаны с недетерминированностью LLM, автоматическими semantic/behavior-oracles и независимой экспертной разметкой, которая ещё не выполнена. Сопоставимый controlled запуск 30B отложен, поэтому преимущество одного размера модели над другим не заявляется. Работа также не утверждает универсальную стабильность, превосходство над всеми альтернативами или полную поддержку Vensim.

Дальнейшая оценка должна включать двух независимых экспертов с adjudication и показателем согласованности, а для проверки влияния размера --- новый сопоставимый 8B/30B-прогон. Исправления найденных failure classes потребуют отдельного regression и нового unseen corpus; подстройка под уже проанализированный blind corpus недопустима.

Таким образом, поставленная инженерная цель достигнута: создана воспроизводимая платформа для исследования ограниченной генерации системно-динамических моделей. Научная проверка преимуществ предложенного метода получила ясный экспериментальный протокол и измеримые критерии, а не подменяется демонстрацией работающего сервиса.
